\section{Intemperismo}

\subsection{Definição}

O \textbf{intemperismo} é o conjunto de processos \textbf{físicos, químicos e biológicos} que provocam a \textbf{desagregação e/ou decomposição das rochas e dos minerais} quando estes ficam expostos às condições existentes na superfície ou próximo à superfície terrestre \parencite{boggs2012,tucker2001}.

As rochas que se formam em profundidade encontram condições de \textbf{pressão, temperatura e composição química} diferentes daquelas existentes na superfície. Quando essas rochas são expostas à atmosfera, à água, aos organismos e às variações de temperatura, muitos de seus minerais deixam de estar em equilíbrio com o novo ambiente e passam por transformações.

De forma geral, o intemperismo pode produzir:

\begin{itemize}
    \item fragmentos de rocha;
    \item grãos minerais;
    \item minerais secundários, como as argilas;
    \item íons dissolvidos na água;
    \item solos e mantos de alteração.
\end{itemize}

Portanto, o intemperismo constitui uma etapa fundamental do \textbf{ciclo das rochas} e também participa da formação dos \textbf{solos e dos sedimentos}, que posteriormente podem ser transportados e depositados.

É importante distinguir \textbf{intemperismo} de \textbf{erosão}:

\begin{itemize}
    \item \textbf{Intemperismo:} alteração ou desagregação da rocha \textbf{no próprio local};
    \item \textbf{Erosão:} remoção e transporte do material produzido pelo intemperismo.
\end{itemize}

Uma rocha pode, portanto, sofrer intemperismo sem que seus produtos sejam imediatamente transportados.

\subsection{Principais tipos de intemperismo}

Tradicionalmente, o intemperismo é dividido em três grandes grupos:

\begin{enumerate}
    \item \textbf{intemperismo físico ou mecânico};
    \item \textbf{intemperismo químico};
    \item \textbf{intemperismo biológico}.
\end{enumerate}

Na natureza, entretanto, esses processos normalmente atuam \textbf{simultaneamente} e podem intensificar uns aos outros.

\subsection{Intemperismo físico ou mecânico}

O \textbf{intemperismo físico} corresponde à fragmentação da rocha em pedaços menores \textbf{sem alteração significativa da composição química dos minerais}.

A rocha original é quebrada, mas os minerais que a formavam permanecem, inicialmente, com praticamente a mesma composição química.

Por exemplo, um bloco de granito pode ser dividido em fragmentos progressivamente menores, mas esses fragmentos ainda podem conter quartzo, feldspato e mica.

A principal consequência do intemperismo físico é o \textbf{aumento da área superficial da rocha}. Isso aumenta a superfície de contato entre os minerais e:

\begin{itemize}
    \item a água;
    \item o oxigênio;
    \item o dióxido de carbono;
    \item os organismos.
\end{itemize}

Consequentemente, o intemperismo físico pode acelerar o \textbf{intemperismo químico}.

\subsubsection{Alívio de pressão ou descompressão}

Muitas rochas se formam em grandes profundidades, submetidas a pressões elevadas devido ao peso das rochas que estão acima delas.

Quando a erosão remove progressivamente esse material superior, a pressão sobre a rocha diminui. Essa redução de pressão permite uma pequena expansão da rocha.

Como consequência, podem surgir \textbf{fraturas aproximadamente paralelas à superfície}, formando estruturas semelhantes a grandes lâminas ou cascas.

Esse processo é frequentemente chamado de \textbf{esfoliação por descompressão} e pode ser observado, por exemplo, em grandes corpos graníticos.

\subsubsection{Variação de temperatura}

As rochas podem sofrer aquecimento durante o dia e resfriamento durante a noite.

Quando aquecidos, os minerais tendem a se expandir; quando resfriados, tendem a se contrair.

Como diferentes minerais apresentam diferentes coeficientes de expansão térmica, ciclos repetidos de aquecimento e resfriamento podem gerar tensões dentro da rocha.

Ao longo de muitos ciclos, essas tensões podem contribuir para a formação e propagação de microfraturas.

Esse processo pode ser importante principalmente em ambientes submetidos a grandes variações térmicas. Entretanto, a simples variação diária de temperatura nem sempre é suficiente para explicar, isoladamente, uma intensa fragmentação das rochas, pois normalmente diversos mecanismos atuam simultaneamente.

\subsubsection{Congelamento da água em fraturas}

A água pode penetrar em pequenas fraturas existentes nas rochas.

Quando a temperatura cai abaixo do ponto de congelamento, essa água congela. A transformação da água líquida em gelo provoca aumento de volume.

Esse aumento exerce pressão sobre as paredes da fratura.

Quando ocorrem repetidos ciclos de

\[
\text{congelamento}
\rightarrow
\text{degelo}
\rightarrow
\text{congelamento},
\]

a fratura pode aumentar progressivamente até provocar a ruptura da rocha.

Esse mecanismo é particularmente importante em regiões frias e montanhosas sujeitas a frequentes ciclos de congelamento e degelo.

\subsubsection{Cristalização de sais}

Em ambientes onde existe água contendo sais dissolvidos, a água pode penetrar nos poros e fraturas das rochas.

Quando essa água evapora, os sais podem cristalizar.

O crescimento dos cristais dentro dos poros pode exercer pressão sobre os minerais e provocar a fragmentação da rocha.

Esse processo é chamado de \textbf{haloclastia} ou \textbf{intemperismo salino}.

Pode ser especialmente importante em:

\begin{itemize}
    \item ambientes costeiros;
    \item regiões áridas;
    \item áreas submetidas à névoa salina;
    \item rochas muito porosas.
\end{itemize}

Em regiões costeiras, os aerossóis marinhos podem introduzir sais nas superfícies rochosas, contribuindo para sua desagregação.

\subsubsection{Ação mecânica de organismos}

Raízes de plantas podem penetrar em pequenas fraturas.

Com o crescimento da raiz, aumenta a pressão exercida sobre as paredes da fratura, favorecendo a separação dos blocos.

Animais que escavam o solo também podem movimentar partículas, abrir canais e expor novos materiais à água e ao ar.

Nesse caso existe uma interação entre \textbf{intemperismo físico e biológico}.

\subsection{Intemperismo químico}

O \textbf{intemperismo químico} envolve reações que modificam a \textbf{composição química e/ou a estrutura mineralógica da rocha}.

Diferentemente do intemperismo físico, no qual ocorre principalmente fragmentação, no intemperismo químico alguns minerais originais podem ser:

\begin{itemize}
    \item dissolvidos;
    \item transformados;
    \item oxidados;
    \item hidratados;
    \item convertidos em novos minerais.
\end{itemize}

A água possui papel fundamental nesses processos.

Em muitos casos, a água funciona como:

\begin{itemize}
    \item solvente;
    \item meio onde ocorrem as reações;
    \item agente de transporte dos íons produzidos.
\end{itemize}

O intemperismo químico tende a ser particularmente intenso em regiões \textbf{quentes e úmidas}, pois temperaturas elevadas aumentam a velocidade de muitas reações químicas e a abundância de água favorece processos como dissolução e hidrólise.

\subsubsection{Dissolução}

A dissolução ocorre quando um mineral entra em contato com a água e seus componentes passam para a solução na forma de íons.

De maneira simplificada:

\[
\text{mineral sólido} + \text{água}
\rightarrow
\text{íons dissolvidos}
\]

Alguns minerais apresentam elevada solubilidade.

Um exemplo é a \textbf{halita}, formada por cloreto de sódio:

\[
\mathrm{NaCl_{(s)}
\rightarrow
Na^+_{(aq)} + Cl^-_{(aq)}}
\]

O mineral deixa de fazer parte da rocha e seus componentes passam para a água.

Esses íons podem posteriormente ser transportados por águas subterrâneas, rios e outros corpos d'água até alcançarem ambientes de deposição, incluindo os oceanos.

\subsubsection{Carbonatação}

A água da chuva contém dióxido de carbono dissolvido.

O dióxido de carbono reage com a água, formando ácido carbônico:

\[
\mathrm{CO_2 + H_2O
\rightleftharpoons
H_2CO_3}
\]

Embora seja um ácido relativamente fraco, o ácido carbônico aumenta a capacidade da água de reagir com determinados minerais.

Esse processo é especialmente importante em rochas carbonáticas, como os \textbf{calcários}.

A calcita, principal mineral de muitos calcários, pode reagir na presença de água e dióxido de carbono:

\[
\mathrm{CaCO_3 + CO_2 + H_2O
\rightarrow
Ca^{2+} + 2HCO_3^-}
\]

Esse processo favorece a dissolução do calcário.

Ao longo de grandes intervalos de tempo, pode contribuir para a formação de relevo cárstico, caracterizado por feições como:

\begin{itemize}
    \item cavernas;
    \item dolinas;
    \item condutos subterrâneos;
    \item sumidouros.
\end{itemize}

A precipitação posterior do carbonato de cálcio no interior de cavernas pode originar estruturas como estalactites e estalagmites.

\subsubsection{Hidrólise}

A \textbf{hidrólise} é um dos processos mais importantes do intemperismo químico dos \textbf{silicatos}, que constituem grande parte da crosta terrestre.

Nesse processo, minerais reagem com água, íons \(H^+\) e outras espécies presentes na solução.

Como resultado, minerais primários podem ser transformados em \textbf{minerais secundários}, especialmente argilominerais.

Um exemplo importante envolve os \textbf{feldspatos}.

De forma simplificada:

\[
\text{feldspato}
+
\text{água}
+
H^+
\rightarrow
\text{argilomineral}
+
\text{íons dissolvidos}
+
\text{sílica}
\]

Um exemplo balanceado é a alteração do feldspato potássico para caulinita:

\[
\mathrm{
2\,KAlSi_3O_8 + 2\,H_2CO_3 + 9\,H_2O
\rightarrow
Al_2Si_2O_5(OH)_4 + 4\,H_4SiO_4 + 2\,K^+ + 2\,HCO_3^-
}
\]

Nessa reação, a caulinita é um mineral secundário; potássio, bicarbonato e sílica dissolvida podem ser removidos pela água.

O feldspato existente originalmente numa rocha magmática pode, portanto, ser progressivamente transformado em minerais de argila.

Esse processo possui grande importância para a formação e evolução dos solos.

\subsubsection{Oxidação}

A \textbf{oxidação} envolve a perda de elétrons por determinados elementos químicos.

No intemperismo, um dos casos mais importantes envolve o \textbf{ferro}.

Minerais contendo ferro ferroso, \(Fe^{2+}\), podem sofrer oxidação, formando ferro férrico, \(Fe^{3+}\).

De maneira simplificada:

\[
Fe^{2+}
\rightarrow
Fe^{3+}
\]

Na presença de água e oxigênio podem ser formados óxidos e hidróxidos de ferro.

Esses produtos frequentemente apresentam cores:

\begin{itemize}
    \item avermelhadas;
    \item amareladas;
    \item alaranjadas;
    \item marrons.
\end{itemize}

Por isso, solos fortemente intemperizados em regiões tropicais frequentemente apresentam colorações avermelhadas ou amareladas.

\subsubsection{Hidratação}

Na \textbf{hidratação}, moléculas de água são incorporadas à estrutura de determinados minerais.

Essa incorporação pode provocar:

\begin{itemize}
    \item aumento de volume;
    \item alteração da estrutura mineral;
    \item redução da resistência do material.
\end{itemize}

Um exemplo é a transformação da anidrita em gipsita:

\[
\mathrm{CaSO_4 + 2H_2O
\rightarrow
CaSO_4\cdot 2H_2O}
\]

Nesse caso, a entrada de água modifica a estrutura cristalina do mineral.

\subsection{Intemperismo biológico}

O \textbf{intemperismo biológico} corresponde à participação dos seres vivos na alteração e desagregação das rochas.

Na prática, a atividade biológica pode provocar tanto efeitos físicos quanto químicos. Por isso, o intemperismo biológico não deve ser entendido como completamente separado dos demais tipos de intemperismo.

\subsubsection{Ação das raízes}

As raízes das plantas podem penetrar em fraturas existentes nas rochas.

À medida que crescem, exercem pressão mecânica, aumentando essas fraturas e favorecendo a fragmentação.

Ao mesmo tempo, raízes e microrganismos associados ao solo podem modificar quimicamente o ambiente.

\subsubsection{Produção de ácidos orgânicos}

A decomposição da matéria orgânica pode produzir diversos ácidos orgânicos.

Esses compostos podem:

\begin{itemize}
    \item reduzir o pH da água;
    \item aumentar a dissolução de minerais;
    \item facilitar a liberação de determinados elementos químicos.
\end{itemize}

Dessa maneira, os organismos podem acelerar o intemperismo químico.

\subsubsection{Líquens}

Os líquens são associações simbióticas que normalmente envolvem fungos e organismos fotossintetizantes.

Eles conseguem colonizar superfícies rochosas aparentemente pouco favoráveis à vida.

Sua presença pode favorecer a alteração das rochas principalmente pela retenção de umidade junto à superfície e pela produção de substâncias capazes de participar das reações químicas de alteração.

Ao longo do tempo, contribuem para a formação das primeiras camadas de material alterado sobre uma rocha exposta.

\subsection{Minerais e resistência ao intemperismo}

Nem todos os minerais apresentam a mesma resistência ao intemperismo.

Isso ocorre porque muitos minerais se formaram sob condições de temperatura e pressão muito diferentes das existentes na superfície terrestre.

De maneira geral, minerais formados em condições de temperatura mais elevada tendem a apresentar menor estabilidade química nas condições superficiais.

Minerais como:

\begin{itemize}
    \item olivina;
    \item piroxênio;
    \item anfibólio;
\end{itemize}

tendem a sofrer alteração química com relativa facilidade.

O \textbf{quartzo}, ao contrário, apresenta elevada resistência química nas condições superficiais.

Por isso, é extremamente comum encontrar grande quantidade de quartzo em:

\begin{itemize}
    \item areias;
    \item praias;
    \item dunas;
    \item depósitos sedimentares;
    \item solos.
\end{itemize}

Essa resistência ajuda a explicar por que muitas areias mineralogicamente maduras apresentam composição predominantemente quartzosa.

\subsubsection{Relação entre as séries de Bowen e Goldich}

A série de reação de Bowen descreve, de forma simplificada, a ordem de cristalização dos principais minerais silicáticos a partir de um magma em resfriamento.

\begin{figure}[H]
  \centering
  \resizebox{0.92\linewidth}{!}{\begin{tikzpicture}[
    x=1cm,y=1cm,
    >=Latex,
    font=\small,
    titulo/.style={
        draw,
        rounded corners=2pt,
        fill=gray!15,
        font=\bfseries,
        minimum width=3.9cm,
        minimum height=0.8cm,
        align=center
    },
    mineralL/.style={
        draw,
        rounded corners=2pt,
        fill=green!18,
        minimum width=3.6cm,
        minimum height=0.8cm,
        align=center
    },
    mineralR/.style={
        draw,
        rounded corners=2pt,
        fill=blue!12,
        text width=3.25cm,
        minimum height=0.95cm,
        align=center
    },
    tardio/.style={
        draw,
        rounded corners=2pt,
        fill=orange!12,
        minimum width=2.35cm,
        minimum height=0.8cm,
        align=center
    },
    notaL/.style={
        align=center,
        text width=3.0cm,
        font=\footnotesize
    },
    notaR/.style={
        align=center,
        text width=3.3cm,
        font=\footnotesize
    }
]

% -------------------------------------------------
% Título geral
% -------------------------------------------------
\node[font=\bfseries\large] at (0.1,6.0) {Série de Reação de Bowen};

% -------------------------------------------------
% Eixo de temperatura
% -------------------------------------------------
\draw[very thick,-{Latex[length=3mm]}] (-5.4,5.5) -- (-5.4,-5.4);
\node[left] at (-5.55,5.1) {alta};
\node[left] at (-5.55,-4.9) {baixa};
\node[below,align=center] at (-5.4,-5.45) {Temperatura\\decrescente};

% -------------------------------------------------
% Cabeçalhos das colunas
% -------------------------------------------------
\node[titulo] (td) at (-1.9,5.1) {Série descontínua};
\node[titulo] (tc) at ( 2.1,5.1) {Série contínua};

% -------------------------------------------------
% Série descontínua
% -------------------------------------------------
\node[mineralL] (ol) at (-1.9,4.0) {Olivina};
\node[mineralL] (px) at (-1.9,2.8) {Piroxênio};
\node[mineralL] (am) at (-1.9,1.6) {Anfibólio};
\node[mineralL] (bt) at (-1.9,0.4) {Biotita};

\draw[very thick,-{Latex[length=3mm]}] (ol.south) -- (px.north);
\draw[very thick,-{Latex[length=3mm]}] (px.south) -- (am.north);
\draw[very thick,-{Latex[length=3mm]}] (am.south) -- (bt.north);

% -------------------------------------------------
% Série contínua
% -------------------------------------------------
\node[mineralR] (pl1) at (2.1,4.0) {Plagioclásio\\cálcico};
\node[mineralR] (pl2) at (2.1,2.8) {Plagioclásio\\cálcico--sódico};
\node[mineralR] (pl3) at (2.1,1.6) {Plagioclásio\\sódico--cálcico};
\node[mineralR] (pl4) at (2.1,0.4) {Plagioclásio\\sódico};

\draw[very thick,-{Latex[length=3mm]}] (pl1.south) -- (pl2.north);
\draw[very thick,-{Latex[length=3mm]}] (pl2.south) -- (pl3.north);
\draw[very thick,-{Latex[length=3mm]}] (pl3.south) -- (pl4.north);

% -------------------------------------------------
% Notas explicativas (deslocadas para não cruzar setas)
% -------------------------------------------------
\node[notaL] at (-3.5,-0.9)
{cada mineral reage com o magma e é \\substituído por \\outro};

\node[notaR] at (4.1,-0.9)
{a composição muda de rica em Ca\\para rica em Na};

% -------------------------------------------------
% Caixa central: minerais tardios
% -------------------------------------------------
\node[titulo, minimum width=5.5cm] (tf) at (0.1,-1.8) {Minerais tardios};

% -------------------------------------------------
% Conexões até a caixa central (desviadas do texto)
% -------------------------------------------------
\coordinate (ljoin) at (-0.95,-1.15);
\coordinate (rjoin) at ( 1.15,-1.15);

\draw[very thick]
    (bt.south) -- ++(0,-0.15) coordinate (btA)
    -- ++(0.95,0) coordinate (btB)
    -- (btB |- ljoin);
\draw[very thick,-{Latex[length=3mm]}]
    (btB |- ljoin) -- (tf.north west);

\draw[very thick]
    (pl4.south) -- ++(0,-0.15) coordinate (plA)
    -- ++(-0.95,0) coordinate (plB)
    -- (plB |- rjoin);
\draw[very thick,-{Latex[length=3mm]}]
    (plB |- rjoin) -- (tf.north east);

% -------------------------------------------------
% Minerais tardios finais
% -------------------------------------------------
\node[tardio] (kf) at (-2.35,-3.45) {K-feldspato};
\node[tardio] (ms) at ( 0.00,-3.45) {Muscovita};
\node[tardio] (qt) at ( 2.35,-3.45) {Quartzo};

% pontos de saída da caixa central
\coordinate (tfL) at ($(tf.south west)!0.18!(tf.south east)$);
\coordinate (tfC) at ($(tf.south west)!0.50!(tf.south east)$);
\coordinate (tfR) at ($(tf.south west)!0.82!(tf.south east)$);

\draw[very thick,-{Latex[length=3mm]}] (tfL) -- (kf.north);
\draw[very thick,-{Latex[length=3mm]}] (tfC) -- (ms.north);
\draw[very thick,-{Latex[length=3mm]}] (tfR) -- (qt.north);

% -------------------------------------------------
% Legendas laterais afastadas do desenho
% -------------------------------------------------
\node[align=left, text width=2.6cm] at (6.0,4.0)
{minerais de\\alta temperatura};

\node[align=left, text width=2.6cm] at (5.0,-3.45)
{minerais de\\baixa temperatura};

\end{tikzpicture}}
  \caption{Esquema simplificado da série de reação de Bowen.}
  \label{fig:bowen}
\end{figure}

Goldich observou que a suscetibilidade dos minerais ígneos ao intemperismo apresenta uma relação aproximadamente inversa com essa ordem de cristalização: minerais formados em temperaturas mais elevadas, como olivina e piroxênio, tendem a ser menos estáveis nas condições superficiais, enquanto o quartzo está entre os mais resistentes \parencite{goldich1938}. Essa relação é uma tendência geral, não uma regra absoluta, pois fraturamento, clima, pH, drenagem e tempo também controlam as taxas de alteração.

\subsection{Fatores que controlam o intemperismo}

A intensidade e o tipo de intemperismo dependem de diversos fatores, entre os quais se destacam:

\begin{itemize}
    \item clima;
    \item composição mineralógica;
    \item estrutura da rocha;
    \item relevo;
    \item presença de água;
    \item atividade biológica;
    \item tempo de exposição.
\end{itemize}

\subsubsection{Clima}

O clima é um dos fatores mais importantes no controle do intemperismo.

Em \textbf{climas quentes e úmidos}, existe elevada disponibilidade de água e as temperaturas mais altas favorecem a velocidade das reações químicas. Nessas condições, o intemperismo químico tende a ser intenso.

Em \textbf{climas frios}, processos associados aos ciclos de congelamento e degelo podem apresentar maior importância.

Em \textbf{climas áridos}, a menor disponibilidade de água limita muitas reações químicas. Entretanto, processos como cristalização de sais e variações térmicas podem apresentar maior importância relativa.

\subsubsection{Composição e estrutura da rocha}

A composição mineralógica determina em grande parte a resistência da rocha.

Uma rocha contendo minerais altamente resistentes pode permanecer exposta durante longo período sem sofrer intensa alteração química.

Por outro lado, minerais instáveis nas condições superficiais podem ser rapidamente transformados.

A estrutura e a textura da rocha também influenciam o intemperismo.

Uma rocha com muitas fraturas, elevada porosidade ou elevada permeabilidade permite maior penetração de água e maior superfície disponível para reações.

\subsubsection{Relevo}

A posição topográfica influencia a permanência dos materiais produzidos pelo intemperismo.

Em encostas muito inclinadas, os produtos de alteração podem ser removidos rapidamente pela gravidade e pela erosão.

Em terrenos de menor declividade, esses materiais podem permanecer por mais tempo sobre a rocha original, permitindo o desenvolvimento de mantos de alteração e solos mais espessos.

\subsubsection{Tempo}

O tempo também é fundamental para o desenvolvimento do intemperismo.

Quanto maior o período de exposição da rocha às condições superficiais, maior tende a ser seu grau de alteração, desde que existam condições adequadas para que os processos ocorram.

Entretanto, essa relação não é simplesmente linear, pois a formação de solos, a drenagem e a remoção dos produtos de alteração modificam continuamente as condições em que o intemperismo ocorre.

\subsection{Intemperismo e formação do solo}

O intemperismo é um dos processos fundamentais para a \textbf{pedogênese}, isto é, a formação e evolução dos solos.

A alteração das rochas produz:

\begin{itemize}
    \item fragmentos minerais;
    \item partículas finas;
    \item argilominerais;
    \item elementos químicos dissolvidos.
\end{itemize}

Esses materiais passam a interagir com matéria orgânica, organismos, água e atmosfera.

De maneira simplificada:

\[
\boxed{
\text{rocha}
\rightarrow
\text{intemperismo}
\rightarrow
\text{material alterado}
\rightarrow
\text{solo}
}
\]

\subsection{Intemperismo, erosão, transporte e sedimentação}

O intemperismo faz parte de uma sequência de processos fundamentais para a dinâmica superficial da Terra.

De maneira simplificada:

\[
\boxed{
\text{rocha}
\xrightarrow{\text{intemperismo}}
\text{material alterado}
\xrightarrow{\text{erosão}}
\text{transporte}
\xrightarrow{\text{deposição}}
\text{sedimento}
}
\]

Depois da deposição, os sedimentos podem sofrer processos de compactação, cimentação e outras transformações diagenéticas, podendo dar origem às \textbf{rochas sedimentares}.

\subsection{Intemperismo e sedimentos}

O intemperismo possui grande importância para a sedimentologia.

O intemperismo físico tende a produzir fragmentos progressivamente menores, enquanto o intemperismo químico pode destruir determinados minerais e formar novos minerais, como as argilas.

Por isso, um sedimento não representa necessariamente apenas pequenos fragmentos da rocha original.

Durante o intemperismo ocorre também uma espécie de \textbf{seleção mineralógica}.

Minerais menos resistentes tendem a ser:

\begin{itemize}
    \item dissolvidos;
    \item alterados;
    \item transformados em minerais secundários.
\end{itemize}

Minerais mais resistentes, como o quartzo, apresentam maior possibilidade de permanecer no sedimento.

Com sucessivos ciclos de

\[
\text{intemperismo}
+
\text{erosão}
+
\text{transporte}
+
\text{deposição},
\]

um sedimento pode tornar-se progressivamente mais \textbf{maduro mineralogicamente}.

\subsection{Intemperismo em ambientes costeiros}

Nos ambientes costeiros, diversos processos de intemperismo podem atuar simultaneamente.

Entre eles destacam-se:

\begin{itemize}
    \item ação da água;
    \item presença de sais;
    \item ciclos de umedecimento e secagem;
    \item atividade biológica;
    \item dissolução de minerais;
    \item oxidação;
    \item cristalização de sais.
\end{itemize}

A \textbf{haloclastia} pode ser especialmente relevante em superfícies rochosas submetidas à névoa e aos aerossóis marinhos.

A água contendo sais penetra nos poros e nas fraturas das rochas. Com a evaporação da água, pode ocorrer a cristalização dos sais.

Ciclos sucessivos de dissolução e cristalização podem favorecer a desagregação da superfície da rocha.

\subsection{Alteração esferoidal}

Um processo observado em algumas rochas é o chamado \textbf{intemperismo esferoidal} ou \textbf{alteração esferoidal}.

Quando um bloco rochoso fraturado sofre intemperismo químico, suas diferentes regiões não são alteradas com a mesma velocidade.

Os cantos apresentam maior superfície exposta em relação ao volume, seguidos pelas arestas e pelas faces.

De maneira simplificada:

\[
\text{cantos}
>
\text{arestas}
>
\text{faces}
\]

em relação à intensidade da alteração.

Com o tempo, blocos inicialmente angulosos podem adquirir formas progressivamente mais arredondadas.

É importante não confundir esse processo com o arredondamento provocado pelo transporte sedimentar.

No intemperismo esferoidal, o arredondamento ocorre \textbf{no próprio local}, como resultado da alteração diferencial da rocha.

\subsection{Perfil de intemperismo}

Quando o intemperismo atua durante longos períodos, pode desenvolver-se um \textbf{perfil de alteração} sobre a rocha.

De maneira simplificada, pode ocorrer a seguinte sequência vertical:

\[
\begin{array}{c}
\text{solo} \\[4pt]
\downarrow \\[4pt]
\text{material fortemente alterado} \\[4pt]
\downarrow \\[4pt]
\text{rocha parcialmente alterada} \\[4pt]
\downarrow \\[4pt]
\text{rocha sã}
\end{array}
\]

A transição entre essas zonas geralmente é gradual.

A espessura do perfil depende de fatores como:

\begin{itemize}
    \item clima;
    \item tempo;
    \item relevo;
    \item tipo de rocha;
    \item drenagem.
\end{itemize}

Em regiões tropicais úmidas, os perfis de alteração podem atingir grandes espessuras.

\subsection{Síntese}

Para fins didáticos, os três tipos principais de intemperismo podem ser resumidos da seguinte forma:

\paragraph{Intemperismo físico}

Produz principalmente a fragmentação da rocha:

\[
\text{rocha}
\rightarrow
\text{fragmentos menores}
\]

sem modificação química significativa dos minerais.

Entre seus principais mecanismos estão:

\begin{itemize}
    \item descompressão;
    \item congelamento e degelo;
    \item cristalização de sais;
    \item variações térmicas;
    \item ação mecânica de organismos.
\end{itemize}

\paragraph{Intemperismo químico}

Produz alterações na composição ou na estrutura dos minerais:

\[
\text{mineral original}
\rightarrow
\text{mineral secundário e/ou íons dissolvidos}
\]

Entre seus principais processos estão:

\begin{itemize}
    \item dissolução;
    \item hidrólise;
    \item carbonatação;
    \item oxidação;
    \item hidratação.
\end{itemize}

\paragraph{Intemperismo biológico}

Resulta da atividade dos organismos e pode provocar tanto efeitos físicos quanto efeitos químicos sobre as rochas.

\subsection{Importância geológica}

Uma maneira de compreender o intemperismo é considerar que uma rocha exposta na superfície terrestre encontra condições ambientais muito diferentes daquelas em que muitos de seus minerais originalmente se formaram.

Alguns minerais permanecem relativamente estáveis. Outros são alterados, dissolvidos ou transformados em novos minerais.

Ao mesmo tempo, a rocha pode ser progressivamente fraturada e desagregada.

O resultado conjunto desses processos fornece grande parte do material que posteriormente será submetido a:

\[
\boxed{
\text{erosão}
\rightarrow
\text{transporte}
\rightarrow
\text{deposição}
\rightarrow
\text{sedimentação}
}
\]

Por isso, o intemperismo constitui uma ligação fundamental entre a \textbf{geologia das rochas}, a \textbf{geomorfologia}, a \textbf{pedologia} e a \textbf{sedimentologia}, além de exercer importante influência sobre os processos que ocorrem nos ambientes continentais, costeiros e marinhos.


% =========================================================
% TRANSPORTE DE SEDIMENTOS
% =========================================================
